\chapter{Minimalne drzewa rozpinające: algorytmy Prima, Kruskala i DSU}
\label{ch:tg-mst}

W rozdziale \ref{ch:tg-bfs-dfs} skonstruowaliśmy drzewa
przeszukiwania, a w rozdziale \ref{ch:tg-shortest-paths}
minimalizowaliśmy koszt pojedynczej drogi.
Tym razem chcemy połączyć \emph{wszystkie} wierzchołki
danej składowej, płacąc jak najmniej za wybrane krawędzie.
Minimalne drzewo rozpinające jest na ogół różne od
drzewa najkrótszych dróg: pierwsze minimalizuje sumę
wag całego drzewa, drugie osobno odległości od źródła.

\section{Zadanie minimalnego drzewa i lasu rozpinającego}
\label{sec:tg-mst-model}

Przez cały rozdział $G=(V,E)$ oznacza \emph{skończony
graf prosty nieskierowany}, $n=|V|$, $m=|E|$,
a $w\colon E\to\mathbb R$ jest dowolną funkcją wag.
W szczególności dopuszczamy wagi ujemne, zerowe i powtarzające
się. Dla $F\subseteq E$ piszemy
$w(F)=\sum_{e\in F}w(e)$; dla $F=\varnothing$ suma wynosi zero.
Pojęcia drzewa rozpinającego oraz przekroju $\delta_G(S)$
zostały zdefiniowane w rozdziale \ref{ch:tg-drzewa}.
Zbiór krawędzi nazywamy \emph{lasem}, jeśli podgraf,
który wyznacza na zbiorze wszystkich wierzchołków,
nie zawiera cyklu.

\begin{definicja}[Minimalne drzewo i minimalny las rozpinający]
\label{def:tg-mst}
Jeżeli $G$ jest niepusty i spójny, \emph{minimalnym drzewem
rozpinającym} (MST, od \emph{minimum spanning tree})
nazywamy drzewo rozpinające $T$ o najmniejszej
możliwej wartości $w(E(T))$.

Dla dowolnego $G$ \emph{minimalnym lasem rozpinającym}
(MSF, od \emph{minimum spanning forest}) nazywamy las
$F\subseteq E$, którego ograniczenie do każdej składowej
spójnej $G$ jest minimalnym drzewem rozpinającym
tej składowej. Wierzchołek izolowany tworzy drzewo
jednowierzchołkowe o koszcie zero.
\end{definicja}

\begin{stwierdzenie}[Istnienie i liczba krawędzi]
\label{prop:tg-mst-existence}
Niech $G=(V,E)$ będzie skończonym grafem prostym
nieskierowanym z wagami $w:E\to\mathbb R$ i $n=|V|$.
Wówczas $G$ ma co najmniej jeden MSF.
Jeżeli $G$ ma $c$ składowych spójnych (dla $V=\varnothing$
przyjmujemy $c=0$), to każdy MSF zawiera dokładnie $n-c$
krawędzi.
\end{stwierdzenie}
\begin{proof}
Każda składowa skończonego grafu ma drzewo rozpinające,
a takich drzew jest skończenie wiele.
Wybierzmy w każdej składowej drzewo o najmniejszej sumie wag.
Ich suma rozłączna daje MSF. Drzewo na $n_i$ wierzchołkach
ma $n_i-1$ krawędzi. Sumując po $c$ składowych,
otrzymujemy $\sum_i(n_i-1)=n-c$.
\end{proof}

\begin{uwaga}[Wagi ujemne nie stanowią problemu]
Algorytmy Prima i Kruskala są poprawne również dla wag
ujemnych. Nie stosują warunku nieujemności używanego
w algorytmie Dijkstry. Nie należy również mylić MSF
z minimalnym podgrafem spójnym o dowolnej liczbie krawędzi:
ujemne wagi mogą sprawić, że dodanie krawędzi
tworzącej cykl obniży koszt. Na przykład w trójkącie
z trzema wagami $-1$ każde drzewo rozpinające ma koszt
$-2$, a cały trójkąt jest spójny i ma koszt $-3$.
\end{uwaga}

\section{Własność przekroju i własność cyklu}
\label{sec:tg-mst-cut-cycle}

\begin{twierdzenie}[Własność przekroju -- bezpieczna krawędź]
\label{thm:tg-mst-cut}
Niech $G=(V,E)$ będzie skończonym grafem prostym
nieskierowanym z wagami $w:E\to\mathbb R$,
a $F\subseteq E$ lasem zawartym w pewnym MSF.
Niech $S\subseteq V$ będzie zbiorem takim, że żadna
krawędź $F$ nie należy do $\delta_G(S)$.
Jeżeli $\delta_G(S)\ne\varnothing$ i
$e\in\delta_G(S)$ jest krawędzią o najmniejszej
wadze w tym przekroju, to istnieje MSF
zawierający $F\cup\{e\}$.
\end{twierdzenie}
\begin{proof}
Wybierzmy MSF $T$ zawierający $F$.
Jeśli $e\in T$, nie ma czego dowodzić.
W przeciwnym przypadku końce $e$ należą
do tej samej składowej spójnej $G$, zatem
dodanie $e$ do drzewa $T$ tej składowej
tworzy dokładnie jeden cykl $C$.
Ponieważ $e$ przecina podział $(S,V\setminus S)$,
ścieżka w $T$ łącząca końce $e$ zaczyna się po jednej
stronie przekroju, a kończy po drugiej. Musi zatem
zawierać co najmniej jedną krawędź przekroju. Wybierzmy
$f\in C\cap\delta_G(S)$, $f\ne e$.
Z założenia $w(e)\leq w(f)$.
Usunięcie $f$ z jedynego cyklu w $T\cup\{e\}$
nie niszczy spójności odpowiedniej składowej,
gdyż końce $f$ nadal łączy pozostała część cyklu.
Usuwa zarazem jedyny cykl. Zatem
$(T\setminus\{f\})\cup\{e\}$ jest lasem rozpinającym
te same składowe, o wadze $w(T)-w(f)+w(e)\leq w(T)$.
Minimalność $T$ wymusza równość kosztów.
Ponadto $f\notin F$, ponieważ przekrój nie
zawiera krawędzi $F$. Otrzymany MSF zawiera
$F\cup\{e\}$.
\end{proof}

\begin{wniosek}[Krawędź jednoznacznie najlżejsza]
\label{cor:tg-mst-strict-cut}
Niech $G=(V,E)$ będzie skończonym grafem prostym
nieskierowanym z wagami $w:E\to\mathbb R$ i $S\subseteq V$.
Jeżeli $e\in\delta_G(S)$ jest jedyną krawędzią o najmniejszej
wadze na tym przekroju,
to $e$ należy do każdego MSF grafu $G$.
\end{wniosek}
\begin{proof}
Gdyby pewien MSF nie zawierał $e$, jego ścieżka
między końcami $e$ musiałaby przekroczyć ten
przekrój inną krawędzią $f$ o wadze większej
od $w(e)$. Dodanie $e$ tworzy jedyny cykl,
zawierający $f$. Usunięcie $f$ zachowuje spójność
składowej i usuwa cykl, jak w dowodzie
Twierdzenia~\ref{thm:tg-mst-cut}.
Otrzymany las rozpinający miałby koszt mniejszy
o $w(f)-w(e)>0$, sprzecznie z minimalnością.
\end{proof}

\begin{twierdzenie}[Własność cyklu]
\label{thm:tg-mst-cycle}
Niech $G=(V,E)$ będzie skończonym grafem prostym
nieskierowanym z wagami $w:E\to\mathbb R$,
a $C$ cyklem prostym w $G$. Jeśli jedna z jego krawędzi
$e$ ma wagę \emph{ściśle większą} od wag wszystkich
pozostałych krawędzi $C$, to $e$ nie należy
do żadnego MSF grafu $G$.
\end{twierdzenie}
\begin{proof}
Gdyby $e$ należała do MSF $T$, usunięcie
jej dzieliłoby odpowiednie drzewo na dwa
niepuste zbiory wierzchołków $S$ oraz $R$.
Ścieżka $C\setminus\{e\}$ łączy końce $e$,
więc zawiera inną krawędź $f$ między $S$ i $R$.
Krawędź $f$ nie należy do $T\setminus\{e\}$,
gdyż łączy dwa różne drzewa tego lasu.
Dodanie $f$ scala te drzewa bez tworzenia cyklu
i nie zmienia pozostałych składowych.
Z założenia $w(f)<w(e)$, więc las rozpinający
$(T\setminus\{e\})\cup\{f\}$ ma wagę
$w(T)-w(e)+w(f)<w(T)$, co jest sprzecznością.
\end{proof}

\begin{twierdzenie}[Kryterium optymalności i jednoznaczności danego MSF]
\label{thm:tg-mst-certificate}
Niech $G=(V,E)$ będzie skończonym grafem
prostym nieskierowanym z wagami
$w:E\to\mathbb R$.
Niech $T\subseteq E$ będzie lasem,
który w każdej składowej spójnej $G$
jest drzewem rozpinającym.
Dla każdej krawędzi $e=\{u,v\}\notin T$
oznaczmy przez $P_T(u,v)$
jedyną drogę łączącą jej końce
w odpowiednim drzewie lasu $T$; dla $e=\{u,v\}$
piszemy również $P_T(e):=P_T(u,v)$.
Maksimum jest określone, gdyż $u\ne v$.
Wówczas:
\begin{enumerate}[label=(\roman*)]
\item $T$ jest minimalnym lasem
rozpinającym wtedy i tylko wtedy,
gdy dla każdej $e\in E\setminus T$
zachodzi
\[
w(e)\geq\max_{f\in P_T(u,v)}w(f).
\]
\item $T$ jest \emph{jedynym}
minimalnym lasem rozpinającym
wtedy i tylko wtedy, gdy wszystkie
powyższe nierówności są ścisłe.
\end{enumerate}
Dla $E\setminus T=\varnothing$
warunki są spełnione w sposób pusty.
\end{twierdzenie}

\begin{proof}
Załóżmy, że $T$ jest MSF.
Jeżeli dla pewnej krawędzi
$e=\{u,v\}\notin T$ zachodziłoby
$w(e)<\max_{f\in P_T(u,v)}w(f)$,
to wybór najcięższej krawędzi
$f$ na tej drodze i wymiana
$T':=(T\setminus\{f\})\cup\{e\}$
dałyby las rozpinający
o ściśle mniejszej wadze.
To sprzeczność.
Nierówności są więc konieczne.

Odwrotnie, przypuśćmy,
że wszystkie nierówności zachodzą.
Spośród wszystkich MSF (jest ich skończenie wiele
i istnieje co najmniej jeden) wybierzmy $S$
maksymalizujący liczbę krawędzi wspólnych z $T$.
Jeżeli $S\ne T$, to ponieważ oba lasy rozpinające
mają tyle samo krawędzi, istnieje $f\in T\setminus S$.
Usunięcie $f$ z $T$ rozdziela jedno z jego drzew
na dwie części. Niech $U$ będzie zbiorem wierzchołków
jednej z nich.
Wówczas $\delta_G(U)\cap T=\{f\}$: pozostałe drzewa
lasu $T$ leżą poza $U$, a w drzewie zawierającym $f$
żadna inna krawędź nie łączy rozdzielonych części.
Końce $f$ są połączone w $S$, więc jedyna ścieżka
między nimi w $S$ przecina ten przekrój pewną
krawędzią $e$. Ponieważ $f\notin S$, mamy
$e\in S\setminus T$. Dodanie $f$ tworzy cykl
zawierający $e$.
Droga $P_T(e)$ musi
zawierać $f$, a więc
$w(e)\geq\max_{g\in P_T(e)}w(g)\geq w(f)$.
Wymiana $S':=(S\setminus\{e\})\cup\{f\}$
nie zwiększa kosztu $S$,
zatem z minimalności $S$
otrzymujemy kolejny MSF.
Ponieważ $e\notin T$ i $f\in T\setminus S$,
mamy $|S'\cap T|=|S\cap T|+1$,
co przeczy wyborowi $S$.
Stąd $T=S$ i $T$ jest MSF.

Jeżeli nierówności są ścisłe,
to w powyższym argumencie
każda taka wymiana z
$S\ne T$ ściśle zmniejszałaby
koszt $S$, co jest niemożliwe.
Zatem żaden inny MSF nie istnieje.
Jeśli natomiast dla pewnej
$e\notin T$ zachodzi równość
$w(e)=\max_{f\in P_T(e)}w(f)$,
wymiana najcięższej krawędzi
$f$ z tej drogi na $e$
zachowuje koszt, ale daje
inny MSF. To dowodzi
warunku koniecznego
jednoznaczności.
\end{proof}

\begin{uwaga}[Sprawdzanie wyniku algorytmu]
Kryterium~\ref{thm:tg-mst-certificate}
stanowi niezależny test poprawności
zwróconego drzewa lub lasu:
wystarczy sprawdzić, że $T\subseteq E$ jest lasem
rozpinającym każdą składową (np. przez DFS,
kontrolując spójność i liczbę krawędzi),
i porównać
wagę każdej krawędzi spoza $T$
z maksymalną wagą na jedynej
drodze drzewowej między jej końcami.
Po przygotowaniu odpowiednich danych
do zapytań o maksimum na drodze
(np. przez podnoszenie binarne)
otrzymujemy procedurę z kosztem
$O((n+m)\log(n+1))$.
Warunek działa również przy
wagach ujemnych, zerowych
i powtarzających się.
\end{uwaga}

\begin{wniosek}[Luka optymalności i odporność jednoznacznego MSF]
\label{cor:tg-mst-stability-gap}
Niech $G=(V,E)$ będzie skończonym
grafem prostym nieskierowanym z funkcją wag
$w:E\to\mathbb R$, a $T\subseteq E$
niech będzie minimalnym lasem rozpinającym.
Dla $e=\{u,v\}\in E\setminus T$ niech $P_T(u,v)$
oznacza zbiór krawędzi jedynej drogi w $T$ między
$u$ i $v$. Połóżmy
\[
\gamma_e(T):=
w(e)-\max_{f\in P_T(u,v)}w(f),
\qquad
\gamma(T):=
\min_{e\in E\setminus T}\gamma_e(T).
\]
Jeżeli $E\setminus T=\varnothing$,
przyjmujemy $\gamma(T)=+\infty$.
Wówczas $\gamma(T)\ge0$, a
$T$ jest jedynym MSF
dokładnie wtedy, gdy
$\gamma(T)>0$.
Ponadto, jeśli $\gamma(T)>0$
i zmienimy każdą wagę o
co najwyżej $\varepsilon\ge0$,
tak że
$|\widetilde w(e)-w(e)|\le\varepsilon$
dla wszystkich $e\in E$
oraz $2\varepsilon<\gamma(T)$,
to $T$ pozostaje
\emph{jedynym} minimalnym
lasem rozpinającym
dla wag $\widetilde w$.
Jeśli $E\setminus T=\varnothing$,
$T$ jest jedynym lasem
rozpinającym niezależnie
od wag.
\end{wniosek}

\begin{proof}
Z kryterium
\ref{thm:tg-mst-certificate}
wszystkie $\gamma_e(T)$
są nieujemne, a ich
ścisła dodatniość
dla wszystkich $e\notin T$
jest równoważna
jednoznaczności. Jeżeli $E\setminus T$ jest
niepusty, minimum jest osiągane, ponieważ zbiór
krawędzi jest skończony. Stąd ścisła dodatniość
wszystkich różnic jest równoważna $\gamma(T)>0$.
Jeżeli $e=uv\notin T$,
to po zmianie wag
\[
\begin{aligned}
\widetilde w(e)-
\max_{f\in P_T(u,v)}\widetilde w(f)
&\ge w(e)-\varepsilon
-\max_{f\in P_T(u,v)}(w(f)+\varepsilon)\\
&=\gamma_e(T)-2\varepsilon\\
&\ge\gamma(T)-2\varepsilon>0.
\end{aligned}
\]
Wszystkie nierówności
certyfikatu optymalności
pozostają więc ścisłe.
Ponieważ zmiana samych wag
nie zmienia spójności
ani zbioru drzew rozpinających,
$T$ nadal jest jedynym MSF.
Przy $E\setminus T=\varnothing$
graf sam jest lasem i
nie ma innego lasu rozpinającego
jego składowe, niezależnie
od przypisanych wag.
\end{proof}

\begin{uwaga}[Ostrość promienia odporności]
Jeżeli $T$ jest jedynym MSF i $\gamma(T)<+\infty$,
to wybieramy $e\notin T$ osiągającą minimum
i $f\in P_T(e)$ osiągającą maksimum w definicji
$\gamma_e(T)$. Zmieniamy wyłącznie wagi
$w(e)$ na $w(e)-\gamma(T)/2$ oraz $w(f)$
na $w(f)+\gamma(T)/2$.
Wtedy $\widetilde w(e)=\widetilde w(f)$.
Każda nierówność certyfikatu
\ref{thm:tg-mst-certificate} pozostaje nieostra,
ponieważ może pogorszyć się najwyżej o $\gamma(T)$.
Stąd $T$ jest nadal MSF, ale wymiana $f$ na $e$
daje inny MSF o tej samej wadze.
Jeżeli dopuścimy zaburzenie o normie maksimum
$\varepsilon>\gamma(T)/2$, zmniejszenie $w(e)$
i zwiększenie $w(f)$ o $\varepsilon$ sprawia,
że wymiana ściśle zmniejsza wagę lasu.
W tym sensie $\gamma(T)/2$ jest dokładnym
promieniem odporności jednoznaczności na
dowolne niezależne zaburzenia wag.
\end{uwaga}

\begin{przyklad}[Granica odporności na perturbacje wag]
\label{ex:tg-mst-stability-gap}
W trójkącie o wagach
$w(ab)=1$, $w(bc)=2$,
$w(ac)=5$ jedynym MST
jest $T=\{ab,bc\}$, przy czym
$\gamma(T)=5-\max\{1,2\}=3$.
Zmiana każdej wagi o co
najwyżej $\varepsilon<3/2$
zachowuje to drzewo jako
jedyny MST. Przy
$\varepsilon=3/2$ możemy
zwiększyć wagę $bc$ do
$7/2$ i zmniejszyć
wagę $ac$ do $7/2$.
Pozostawiając wagę $ab$ równą $1$, otrzymujemy dwa
drzewa $\{ab,bc\}$ i $\{ab,ac\}$ o koszcie $9/2$.
Trzecie drzewo $\{bc,ac\}$ ma koszt $7$.
Pierwsze dwa drzewa są zatem minimalne, więc
na granicy oszacowania jednoznaczność może zaniknąć.
\end{przyklad}

\begin{stwierdzenie}[Jednoznaczność przy różnych wagach]
\label{prop:tg-mst-unique}
Jeżeli wagi wszystkich krawędzi skończonego
grafu $G$ są parami różne, to $G$ ma dokładnie
jeden MSF. Odwrotna implikacja nie zachodzi.
\end{stwierdzenie}
\begin{proof}
Przypuśćmy, że istnieją dwa różne MSF $T_1,T_2$.
Wybierzmy krawędź $e$ o najmniejszej wadze
w różnicy symetrycznej ich zbiorów krawędzi
i, po ewentualnej zamianie oznaczeń, przyjmijmy
$e\in T_1\setminus T_2$.
Dodanie $e$ do $T_2$ tworzy cykl. Pewna krawędź
$f$ tego cyklu należy do $T_2\setminus T_1$:
gdyby wszystkie pozostałe krawędzie cyklu
należały do $T_1$, wraz z $e$ tworzyłyby
cykl w $T_1$. Z wyboru $e$ i różności wag
mamy $w(e)<w(f)$, więc
$(T_2\setminus\{f\})\cup\{e\}$
ma koszt mniejszy niż $T_2$, sprzeczność.
Jako kontrprzykład do implikacji odwrotnej
wystarczy graf będący drzewem o dwóch
krawędziach o tej samej wadze: samo to drzewo
jest jedynym MST.
\end{proof}

\section{Kruskal: porządkowanie krawędzi}
\label{sec:tg-mst-kruskal}

Algorytm Kruskala porządkuje wszystkie krawędzie
według niemalejących wag. Przy równych wagach
stosujemy dowolny, lecz ustalony porządek
rozstrzygania remisów. Zaczynamy od
$F=\varnothing$ i przeglądamy kolejne krawędzie
$e=\{u,v\}$. Dodajemy $e$ dokładnie wtedy,
gdy $u$ i $v$ należą do różnych składowych
lasu $(V,F)$. Krawędzie tworzące cykl pomijamy.

\begin{twierdzenie}[Poprawność Kruskala]
\label{thm:tg-mst-kruskal}
Dla każdego skończonego nieskierowanego
grafu ważonego algorytm Kruskala zwraca MSF,
także jeśli graf jest niespójny albo
ma ujemne bądź powtarzające się wagi.
\end{twierdzenie}
\begin{proof}
W każdym kroku $F$ jest lasem, bo dodajemy
krawędź tylko między różnymi jego składowymi.
Dowodzimy indukcyjnie względem liczby rozpatrzonych
krawędzi, że $F$ zawiera się w pewnym MSF.
Dla $F=\varnothing$ teza wynika z
\ref{prop:tg-mst-existence}.
Jeśli kolejna krawędź ma oba końce w tej samej
składowej $(V,F)$, zostaje pominięta i niezmiennik
pozostaje prawdziwy. Niech zatem rozpatrywana
krawędź $e=uv$ łączy dwa różne drzewa lasu.
Weźmy za $S$ zbiór wierzchołków
drzewa $F$ zawierającego $u$.
Żadna krawędź $F$ nie przecina
$\delta_G(S)$. Gdyby istniała
krawędź $f$ tego przekroju
o wadze mniejszej niż $w(e)$,
zostałaby rozpatrzona wcześniej.
Zbiór wybranych krawędzi może jedynie rosnąć,
więc składowe lasu mogą się łączyć, ale nie rozdzielać.
Skoro końce $f$ są teraz rozdzielone, były też
rozdzielone w chwili rozpatrywania $f$.
Algorytm musiałby zatem dodać $f$,
sprzecznie z definicją $S$.
Krawędź $e$ jest więc
jedną z najlżejszych na tym
przekroju, a twierdzenie
\ref{thm:tg-mst-cut} przedłuża
indukcyjny niezmiennik.

Po przejrzeniu całej listy
żadne dwa wierzchołki połączone
ścieżką w $G$ nie mogą leżeć
w różnych składowych $(V,F)$:
na takiej ścieżce znalazłaby
się krawędź między tymi
składowymi, którą algorytm
musiałby dodać podczas
przetwarzania. Stąd $F$
rozpina każdą składową $G$.
Ponieważ $F$ jest lasem rozpinającym, ma $n-c$
krawędzi, gdzie $c$ jest liczbą składowych $G$.
Z niezmiennika $F\subseteq T$ dla pewnego MSF $T$
oraz $|T|=n-c=|F|$ wynika $F=T$.
\end{proof}

\section{Struktura DSU: zbiory rozłączne}
\label{sec:tg-mst-dsu}

W Kruskalu trzeba szybko sprawdzać,
czy końce krawędzi są już połączone.
Służy do tego DSU (ang. \emph{disjoint-set union}),
zwane również strukturą \emph{union--find}.

\begin{definicja}[Operacje DSU]
\label{def:tg-mst-dsu}
Dla skończonego zbioru $V$ przechowujemy podział
na rozłączne klasy. Operacja
$\mathrm{find}(v)$ zwraca reprezentanta
klasy zawierającej $v$.
Operacja $\mathrm{union}(u,v)$
scala klasy zawierające $u$ i $v$,
jeśli są różne.
Początkowo każda klasa jest
jednoelementowa.
\end{definicja}

Realizacja tablicowa przechowuje dla
każdego $v$ rodzica $\mathrm{parent}[v]$
oraz, dla korzeni, rozmiar $\mathrm{size}[v]$.
W operacji $\mathrm{find}$ podążamy
po rodzicach aż do korzenia
i opcjonalnie skracamy ścieżkę,
przypisując wszystkim odwiedzonym
wierzchołkom tego samego korzenia.
Przy łączeniu dwóch różnych klas
korzeń klasy o mniejszym rozmiarze zostaje
dzieckiem korzenia klasy o większym rozmiarze
(\emph{union by size}); przy równości rozmiarów
wybieramy dowolny z dwóch korzeni. W nowym korzeniu
zapisujemy sumę rozmiarów obu klas. Początkowo
$\mathrm{parent}[v]=v$ oraz $\mathrm{size}[v]=1$.
Korzeń rozpoznajemy po warunku $\mathrm{parent}[v]=v$.

\begin{stwierdzenie}[Poprawność i koszt DSU]
\label{prop:tg-mst-dsu}
Operacje $\mathrm{find}$ i $\mathrm{union}$
zachowują dokładnie podział uzyskany
przez dotychczasowe scalenia.
Przy łączeniu według rozmiaru
wysokość drzew rodziców wynosi
co najwyżej $\lfloor\log_2 n\rfloor$
dla $n\geq1$. Dlatego nawet
bez kompresji ścieżek każda
operacja kosztuje $O(\log(n+1))$.
Inicjalizacja wymaga $O(n)$ pamięci
i czasu.
\end{stwierdzenie}
\begin{proof}
Na początku każdy wierzchołek
jest korzeniem swojej klasy.
Jeśli oba argumenty operacji $\mathrm{union}$
mają ten sam korzeń, podział się nie zmienia.
W przeciwnym przypadku zmieniamy rodzica jednego
z dwóch korzeni na drugi korzeń. Drzewa były rozłączne,
więc nie mogła istnieć droga rodziców z drugiego
korzenia do pierwszego. Nowe połączenie nie tworzy
cyklu (poza konwencjonalną pętlą rodzica korzenia)
i wszystkie wierzchołki obu klas mają teraz wspólny
korzeń. Korzenie pozostałych klas nie zmieniają się.
Kompresja ścieżek nie
zmienia korzenia żadnego
wierzchołka, zatem nie
zmienia podziału. Początkowe rozmiary są poprawne,
a przy scalaniu rozłącznych klas ich suma jest
rozmiarem nowej klasy. Kompresja nie zmienia
liczby elementów w żadnej klasie. Dane o rozmiarach
są więc poprawne w każdym korzeniu.

Jeśli głębokość danego wierzchołka
zwiększa się o jeden podczas
łączenia według rozmiaru, to
liczba elementów w jego klasie
co najmniej się podwaja.
Ponieważ rozmiar klasy
nie przekracza $n$,
głębokość może wzrosnąć
co najwyżej $\lfloor\log_2 n\rfloor$ razy.
Kompresja ścieżek nie zwiększa głębokości żadnego
wierzchołka: zmieniane wskaźniki prowadzą wprost
do dotychczasowego korzenia. Głębokość początkowa
wynosi zero, a zatem bieżąca głębokość nie przekracza
liczby dotychczasowych podwojeń rozmiaru klasy.
Operacja $\mathrm{find}$ przechodzi wzdłuż co najwyżej
tylu krawędzi drzewa; ewentualna kompresja wymaga
co najwyżej liniowej pracy względem długości tej drogi.
Operacja $\mathrm{union}$ wykonuje dwa wyszukania
korzenia i stałą liczbę aktualizacji tablic.
\end{proof}

\begin{uwaga}[Silniejsze oszacowanie amortyzowane]
Dla połączenia kompresji ścieżek
z łączeniem według rozmiaru
znane klasyczne twierdzenie
o DSU daje amortyzowany koszt
$O(\alpha(n))$ jednej operacji,
gdzie $\alpha$ jest odwrotnością
funkcji Ackermanna w odpowiedniej
standardowej normalizacji.
Pełny dowód tego mocniejszego
oszacowania wymaga osobnej
analizy amortyzowanej i nie
jest tu przedstawiany.
Udowodnione wyżej ograniczenie
logarytmiczne wystarcza
do oszacowania Kruskala.
\end{uwaga}

\begin{stwierdzenie}[Złożoność Kruskala z DSU]
\label{prop:tg-mst-kruskal-cost}
Dla grafu o $n$ wierzchołkach
i $m$ krawędziach Kruskal
z sortowaniem porównaniowym
i DSU przez łączenie według
rozmiaru działa w czasie
\[
O\bigl(n+m\log(m+1)+
      m\log(n+1)\bigr)
\]
i używa $O(n+m)$ pamięci.
Z kompresją ścieżek można
dodatkowo wykorzystać
wspomniane oszacowanie
amortyzowane $O(\alpha(n))$
dla operacji DSU.
\end{stwierdzenie}
\begin{proof}
Inicjalizacja DSU kosztuje $O(n)$,
sortowanie $O(m\log(m+1))$.
Dla każdej krawędzi wykonujemy
stałą liczbę operacji
$\mathrm{find}$ i co najwyżej
jedną operację $\mathrm{union}$,
o koszcie $O(\log(n+1))$ każda.
Lista posortowanych krawędzi
i tablice DSU zajmują $O(n+m)$
pamięci.
\end{proof}

\section{Prim: rozwijanie jednego drzewa}
\label{sec:tg-mst-prim}

Algorytm Prima dla spójnego
grafu rozpoczynamy od korzenia
$r$ i zbioru $S=\{r\}$.
Początkowo wstawiamy do kopca wszystkie krawędzie
incydentne z $r$. Kopiec jest uporządkowany według wag
i zawiera wszystkie krawędzie mające jeden koniec
w $S$, a drugi poza $S$; może również zawierać
nieaktualne wpisy krawędzi o obu końcach w $S$.
Pobieramy najlżejszą z nich,
dołączamy nowy wierzchołek
i dodajemy do kopca jego
krawędzie prowadzące
do jeszcze nieodwiedzonych
wierzchołków. Wpisy,
których oba końce zostały
już odwiedzone, odrzucamy.
Jeżeli kolejka się opróżni,
zakończyliśmy tę składową.
W grafie niespójnym zaczynamy
od kolejnego nieodwiedzonego
wierzchołka i powtarzamy
procedurę, otrzymując MSF.

\begin{twierdzenie}[Poprawność Prima]
\label{thm:tg-mst-prim}
Dla dowolnego skończonego
nieskierowanego grafu ważonego
algorytm Prima, uruchomiony
kolejno na wszystkich
składowych spójnych,
zwraca minimalny las
rozpinający.
\end{twierdzenie}
\begin{proof}
W obrębie jednej składowej
utrzymujemy drzewo $F$
rozpinające zbiór
już odwiedzonych wierzchołków
$S$. Dla $S=\{r\}$ jest ono
puste. Każda wybierana
krawędź łączy $S$ z
$V\setminus S$, więc
nie tworzy cyklu.
Każda krawędź przekroju została wstawiona do kopca
przy odwiedzeniu jej końca należącego do $S$.
Nie mogła zostać wcześniej usunięta: pobranie aktualnej
krawędzi dołącza jej drugi koniec do $S$, a odrzucenie
nieaktualnej wymaga, aby oba końce już należały do $S$.
Zatem po usunięciu wpisów nieaktualnych
najtańszy jego element
jest najlżejszą krawędzią
przekroju $\delta_G(S)$.
Żadna krawędź $F$
nie przecina tego podziału,
więc twierdzenie
\ref{thm:tg-mst-cut} zapewnia,
że dodanie wybranej
krawędzi zachowuje
możliwość przedłużenia
$F$ do MSF.
Gdy kopiec jest pusty, z $S$ nie wychodzi żadna
krawędź. Gdyby w tej samej składowej istniał
wierzchołek poza $S$, ścieżka od $r$ do niego
musiałaby przeciąć przekrój $\delta_G(S)$.
Zatem $S$ jest całą aktualną składową spójną.
Drzewo $F$ rozpina tę składową i zawiera się
w pewnym MSF, więc jest minimalnym drzewem
rozpinającym tej składowej. Stosując argument
do każdej składowej, otrzymujemy MSF; dla grafu
pustego wynikiem jest pusty las.
\end{proof}

\begin{stwierdzenie}[Złożoność Prima]
\label{prop:tg-mst-prim-cost}
Wariant Prima z listami
sąsiedztwa i kopcem
krawędzi, dopuszczającym
wpisy nieaktualne,
wykonuje $O(n+m\log(m+1))$
operacji i potrzebuje
$O(n+m)$ pamięci
łącznie z kopiami
przechowywanych krawędzi.
\end{stwierdzenie}
\begin{proof}
Każdy wierzchołek
odwiedzamy raz.
Każda krawędź zostaje
wstawiona do kopca
co najwyżej raz:
podczas odwiedzenia
wcześniejszego z jej
końców, jeśli drugi
jest wówczas nieodwiedzony.
Liczba operacji kopcowych
jest więc $O(m)$,
a rozmiar kopca nie
przekracza $m$.
Każde wstawienie i pobranie kosztuje
$O(\log(m+1))$ dla $m\geq1$. Odczyty list
sąsiedztwa kosztują razem $O(n+m)$, gdyż
każda krawędź nieskierowana występuje na dwóch listach.
Tablica odwiedzonych wierzchołków i przejście po
potencjalnych korzeniach kosztują $O(n)$ czasu,
również dla $m=0$. Kopiec zajmuje $O(m)$ pamięci,
a tablice i wynikowy las $O(n)$.
\end{proof}

\section{Przykłady, remisy i kontrola wyników}
\label{sec:tg-mst-examples}

\begin{przyklad}[Ten sam koszt, inne kolejności]
\label{ex:tg-mst-five}
Niech $V=\{a,b,c,d,e\}$,
a krawędzie mają wagi
$w(ab)=1$, $w(bc)=2$,
$w(cd)=3$, $w(de)=2$,
$w(ae)=9$, $w(ac)=4$,
$w(bd)=5$, $w(ce)=6$.
Kruskal wybiera $ab$,
następnie $bc$ i $de$
(w dowolnej kolejności
dla dwóch wag $2$),
a w końcu $cd$.
Otrzymuje drzewo
o koszcie $1+2+2+3=8$.
Prim z korzeniem $a$
może wybierać kolejno
$ab,bc,cd,de$, również
otrzymując koszt $8$.
\end{przyklad}

\begin{figure}[htbp]
\centering
\begin{tikzpicture}[scale=.95]
\node[tgBlueVertex] (a) at (-2,0) {$a$};
\node[tgVertex] (b) at (-1,1.5) {$b$};
\node[tgVertex] (c) at (1,1.5) {$c$};
\node[tgVertex] (d) at (2,0) {$d$};
\node[tgVertex] (e) at (0,-1.7) {$e$};
\draw[tgGhost] (a) to[bend right=12] node[below left,font=\scriptsize] {$9$} (e);
\draw[tgGhost] (a) to[bend right=20]
  node[pos=.48,below,font=\scriptsize,fill=white,inner sep=1pt] {$4$} (c);
\draw[tgGhost] (b) .. controls (.0,2.85) and (3.05,2.65) ..
  node[pos=.53,above,font=\scriptsize,fill=white,inner sep=1pt] {$5$} (d);
\draw[tgGhost] (c)--node[right,font=\scriptsize] {$6$} (e);
\draw[tgTree] (a)--node[above left,font=\scriptsize] {$1$} (b);
\draw[tgTree] (b)--node[above,font=\scriptsize] {$2$} (c);
\draw[tgTree] (c)--node[above right,font=\scriptsize] {$3$} (d);
\draw[tgTree] (d)--node[below right,font=\scriptsize] {$2$} (e);
\end{tikzpicture}
\caption{Minimalne drzewo rozpinające o koszcie $8$:
krawędzie wybrane zaznaczono linią ciągłą,
a pominięte -- przerywaną.}
\label{fig:tg-mst-five}
\end{figure}

\begin{przyklad}[Remisy nie muszą dawać jednoznaczności]
\label{ex:tg-mst-ties}
Dla trójkąta z wagami wszystkich
trzech krawędzi równymi $1$
istnieją dokładnie trzy różne MST,
każde o koszcie $2$.
Jeśli jednak graf jest drzewem,
jego jedynym MST jest to drzewo
nawet przy powtarzających się wagach.
\end{przyklad}

\begin{przyklad}[Graf niespójny i waga ujemna]
\label{ex:tg-mst-disconnected}
Niech $G$ ma składową
o wierzchołkach $\{a,b,c\}$
z wagami $w(ab)=-3$,
$w(bc)=1$, $w(ac)=4$,
drugą składową będącą
krawędzią $de$ o wadze $0$,
oraz izolowany wierzchołek $f$.
MSF zawiera $ab$, $bc$ i $de$,
ma $6-3=3$ krawędzie
i koszt $-3+1+0=-2$.
Nie istnieje drzewo
rozpinające cały $G$,
bo $G$ jest niespójny.
\end{przyklad}

\begin{uwaga}[Dlaczego własność przekroju łączy oba algorytmy]
Kruskal wybiera najlżejszą
krawędź przekroju jednej
z aktualnych składowych lasu,
a Prim -- najlżejszą krawędź
przekroju zbioru już
odwiedzonych wierzchołków.
Różnią się kolejnością,
lecz podstawą poprawności
obu metod jest to samo
twierdzenie \ref{thm:tg-mst-cut}.
W następnym module
własności przekroju
pojawią się ponownie
przy przepływach
i minimalnych przekrojach,
ale tam optymalizujemy
inną wielkość.
\end{uwaga}
